\section{The problem exposed by the repository}
\label{sec:scope}

The root arithmetic README makes a decisive distinction. The
\code{ConciseIPRA} library supplies kernel-checked equivalences between
presentations of intuitionistic primitive recursive arithmetic. The
\code{docs/proofs} directory contains 60 delivered research packets,
including nine framework packets, whose claims were filed as delivered;
the index does not certify them. In particular, the root README says that
none of those packets supplies a completed derivation of its main
application in \code{Prf PRAAx}~\cite{repo-root,repo-index}.
This report uses that distinction throughout. An excellent mathematical
outline, a total Lean function, a successful exact computation, and an
embedded arithmetic proof are four different forms of evidence.

\subsection{What can actually be improved}

The earlier frameworks already supply bounded Booleanization, finite
contracts, precision budgets, conditional inverses, finite nets,
metastability, and several kinds of coherent approximation.
The finite-demand packet treats nonmonotone backward precision schedules
and aggregate root supports. The coherent-methods packet supplies filtered
lifting, Hurwitz coefficients, and weighted determinant norms. The
finite-observation packet contains degree-uniform approximate
factorization, convexity-based continuity bounds, and finite-net
minimization~\cite{repo-demand,repo-coherent,repo-observation}.
Repeating these methods under a new name would not resolve a new difficulty.

The present extension concentrates on the following three gaps.

\begin{enumerate}
\item \textbf{A missing executable syntax bridge.}
The earlier \code{BackwardDemand.lean} defines a semantic
\code{compiledRecursion} and proves its recurrence equations, but explicitly
leaves the construction of its \code{PR} syntax and the object-theory
compiler unimplemented. We supply actual syntax for that program, checked
semantic correctness of the generated syntax, and a general checked
term-to-symbol correctness theorem inside \code{Prf PRAAx}. The remaining
internal proofs are identified individually.
\item \textbf{Termination arguments with resets.}
A single decreasing natural number is often too restrictive for an
algorithm that finishes one task, lowers a more significant counter, and
restarts a lower counter. A fixed finite lexicographic rank suffices. We
derive its execution bound from a primitive-recursive finite simulation,
so a separate numerical runtime bound need not be guessed. The distinction
between a fixed number of counter levels and an input-dependent number is
essential.
\item \textbf{General plane intersections and auxiliary equations.}
The earlier zero-count repairs use monic plane norms, hyperelliptic
curves, or particular weighted leading forms. A bounded shear removes the
coordinate restriction. A Sylvester determinant supplies both the sharp
degree-product count and a complete rational-point procedure for every
intersection passing its nonvanishing test. A finite coefficient envelope
then upgrades pointwise auxiliary witnesses into a uniform complete list.
\end{enumerate}

These contributions extend standard ideas in primitive recursion,
termination analysis, and elimination theory; no claim of historical
priority is made. The gain is an explicit constructive formulation,
quantifier and resource accounting, an executable syntax bridge, and a
precise comparison with the inspected repository proposals.

\subsection{A targeted comparison}

\begin{longtable}{>{\raggedright\arraybackslash}p{0.22\textwidth}>{\raggedright\arraybackslash}p{0.33\textwidth}>{\raggedright\arraybackslash}p{0.37\textwidth}}
\toprule
Existing packet & Reused achievement & Addition or retained boundary\\
\midrule\endhead
\code{ipra_embedding_framework}
& Finite observations; simple analytic branches; hyperelliptic zero counts
& Arbitrary plane intersections passing a resultant test; complete rational lists.\\
\code{ipra_embedding_framework_2026-11-15}
& Finite traces, guarded operations, finite nets, height/list distinction
& A resource hierarchy that constructs fuel for fixed-depth ranked runs.\\
\code{ipra_framework}
& Finite contracts; explicit power-gap families; object-theory infrastructure
& No repetition of the power-gap application; polynomial elimination supplies a different finite-list producer.\\
\code{ipra_framework_2026-11-15}
& Certificates, moduli, elementary bounded descent
& Resetting counters and fixed-versus-variable depth are made explicit.\\
\code{ipra_finite_demand}
& Backward demands; semantic parameter-changing recursion; root supports
& Actual \code{PR} syntax for its recursion. Exact singular Puiseux coefficients and multiplicity matching remain outside this report.\\
\code{ipra_finite_observation}
& Variable-depth tables; complete approximate factorization; restricted curve section bounds
& Finite auxiliary equations yield effective point lists when their coefficient heights are bounded.\\
\code{ipra_coherent_methods_2026-11-15}
& Filtered lifting, coherent towers, weighted norms
& No weight or smoothness assumption in the plane-intersection theorem.\\
\code{ipra_evidence_framework}
& Relational finite outputs and finite compatibility
& Finite exhaustive envelopes remove a pointwise auxiliary choice; no compactness is inferred.\\
\code{ipra_uniform_certificates}
& Uniform finite-net and root certificates
& A syntax-level bridge and a separate termination resource layer.\\
\bottomrule
\end{longtable}

The other four framework notes are indexed together in~\cite{repo-other};
the earlier embedding and contract packages are also consulted directly
\cite{repo-embedding,repo-framework}.
This is a targeted comparison of framework sources and selected difficulty
reports. It is not a proof audit of all 60 application packets.
The provenance inventory records exact paths and hashes for the inspected
sources. Where a prior result is needed, its status remains visible; a
README's self-description is not silently promoted to a verified theorem.

\subsection{The object theory and the ambient metatheory}

We use the standard first-order presentation of \(\ipra\): intuitionistic
first-order logic with equality; symbols for primitive-recursive
definitions, including composition and primitive recursion, with their
defining equations; successor injectivity and separation from zero; and
induction for quantifier-free formulas. Quantifier-free includes the
intuitionistic propositional connectives. Bounded formulas below are
admissible abbreviations after the Booleanization proof in
Section~\ref{sec:framework}.

The metatheory is ordinary mathematics, and Lean is used to check selected
metatheorems and to construct selected object proofs. Thus the following
statements must not be confused:
\[
  \text{Lean proves }\forall n\,P(n),
  \qquad
  \text{Lean constructs }\code{Prf PRAAx []}\,\ulcorner\forall n\,P(n)\urcorner.
\]
Only the second statement directly establishes derivability in the
embedded arithmetic theory. A semantic proof can help discover a finite
arithmetic proof, but no general reflection principle is available to
convert it automatically.

Similarly, \(\forall f\) in a theorem about all primitive-recursive
functions normally belongs to the metatheory: for each fixed definition
of \(f\), construct an arithmetic symbol and proof. It is not a
first-order quantification over an internal type of total functions.
Numerical parameters of a fixed definition, matrix sizes, polynomial
degrees, and lengths of finite lists may vary inside arithmetic.
The permitted variation must be established by the program construction,
not inferred merely from the notation.
